\documentclass[11pt,letterpaper,sans]{moderncv}

% moderncv themes
\moderncvstyle{banking} % style options: 'casual', 'classic', 'banking', 'oldstyle', 'fancy'
\moderncvcolor{blue} % color options: 'black', 'blue', 'burgundy', 'green', 'grey', 'orange', 'purple', 'red'

% character encoding
\usepackage[utf8]{inputenc}
\usepackage[scale=0.82]{geometry}
\usepackage{fontawesome}
\usepackage{academicons}

% Adjust margins and spacing
\geometry{left=1.5cm, right=1.5cm, top=1.2cm, bottom=1.2cm}
\setlength{\hintscolumnwidth}{2.8cm}

% Redefine name color to be consistent
\definecolor{namecolor}{rgb}{0,0,0} % Black color for name

% Personal data
\firstname{\color{namecolor}Anirban}
\familyname{\color{namecolor}Sarkar}
\address{\small Computational Postdoctoral Fellow, Cold Spring Harbor Laboratory, USA}
\phone[mobile]{+1~617~821~5834 / +91~7675~909695}
% \email{asarkar@cshl.edu / anirbansarkar.cs@gmail.com}
\extrainfo{\faEnvelopeO~asarkar@cshl.edu / anirbansarkar.cs@gmail.com}
\homepage{anirbansarkar-cs.github.io}
% \social[github]{anirbansarkar-cs}
% \social[linkedin]{anirban-sarkar}

\begin{document}
\makecvtitle

\vspace{-0.3cm}

% ==================== RESEARCH INTERESTS ====================
\section{Research Interests}
\cvitem{Core Areas}{Machine Learning, Deep Learning, Computer Vision, Computational Biology}
\cvitem{Focus}{Deep Generative Modeling, Explainable AI, Inference-time Control, Trustworthy \& Robust Machine Learning}
\cvitem{Applications}{Model interpretability and robustness under distribution shifts; deep generative models for biological sequence design, inference-time optimization, and evaluation; emerging directions in concept-aligned, cross-modal counterfactual generation}

\vspace{0.2cm}

% ==================== ACADEMIC POSITIONS ====================
\section{Academic Positions}

\cventry{2023--Present}{Computational Postdoctoral Fellow}{Cold Spring Harbor Laboratory}{Cold Spring Harbor, NY, USA}{}{
\textbf{Simons Center for Quantitative Biology} $\cdot$ Advisor: \href{https://www.cshl.edu/research/faculty-staff/peter-koo/}{\textcolor{blue}{Peter Koo}} (\href{https://koolab.cshl.edu/}{\textcolor{blue}{Koo Lab}})\\
\textit{Focus:} Developing deep generative models tailored with biological priors and interpreting them to understand how they build representations to push the boundaries of basic science and cancer biology
}

\vspace{0.2cm}

\cventry{2022--2023}{Postdoctoral Associate}{Massachusetts Institute of Technology}{Cambridge, MA, USA}{}{
\textbf{Department of Brain and Cognitive Sciences} $\cdot$ Advisors: \href{http://www.mit.edu/~xboix/}{\textcolor{blue}{Xavier Boix}} and \href{https://bcs.mit.edu/directory/pawan-sinha}{\textcolor{blue}{Pawan Sinha}} (\href{https://www.sinhalab.mit.edu/}{\textcolor{blue}{Sinha Lab}})\\
\textit{Focus:} Bridging human and machine intelligence through neuroscience-inspired approaches to improve interpretability and address the lack of robustness outside training distributions
}

\vspace{0.2cm}

% ==================== EDUCATION ====================
\section{Education}

\cventry{2016--2022}{Ph.D. in Computer Science}{Indian Institute of Technology Hyderabad}{India}{}{
Advisor: \href{https://people.iith.ac.in/vineethnb/}{\textcolor{blue}{Vineeth N Balasubramanian}}\\
\textit{Research Focus:} Gradient-based and causal explainability techniques, self-explaining neural networks, adversarial and attributional robustness\\
\textbf{Awards:} \href{https://ikdd.acm.org/awards-2022.php}{\textcolor{blue}{IKDD Best Doctoral Dissertation in Data Science Award}} (Winner)
}

\cventry{2014--2016}{M.Tech in Computer Science}{National Institute of Technology Rourkela}{India}{}{
Specialization: Information Security
}

\cventry{2007--2010}{Master of Computer Applications (MCA)}{IIEST Shibpur}{India}{}{}

\vspace{0.2cm}

% ==================== SELECTED PUBLICATIONS ====================
\section{Selected Publications}
\cvitem{}{{\small \textit{Google Scholar Citations: 5400+} $\cdot$ \textit{h-index: 8}}}

\vspace{0.2cm}

\subsection{Recent High-Impact Papers}
\vspace{0.2cm}

\cvitem{MLCB 2024}{\textbf{A. Sarkar}, Z. Tang, C. Zhao, P. Koo. ``\textbf{Designing DNA With Tunable Regulatory Activity Using Discrete Diffusion}''. Machine Learning in Computational Biology. \textit{Long Oral Presentation} $\cdot$ \textit{[41 citations]}}

\vspace{0.1cm}

\cvitem{CVPR 2022}{\textbf{A. Sarkar}, D. Vijaykeerthy, A. Sarkar, VN Balasubramanian. ``\textbf{A Framework for Learning Ante-hoc Explainable Models via Concepts}''. IEEE Conference on Computer Vision and Pattern Recognition. \textit{[104 citations]}}

\vspace{0.1cm}

\cvitem{WACV 2022}{\textbf{A. Sarkar}, A. Sarkar, VN Balasubramanian. ``\textbf{Leveraging Test-Time Consensus Prediction for Robustness against Unseen Noise}''. IEEE Winter Conference on Applications of Computer Vision. \textit{[8 citations]}}

\vspace{0.1cm}

\cvitem{NeurIPS 2021}{\textbf{A. Sarkar}, A. Sarkar, S. Gali, VN Balasubramanian. ``\textbf{Adversarial Robustness without Adversarial Training: A Teacher-Guided Curriculum Learning Approach}''. \textit{[16 citations]}}

\vspace{0.1cm}

\cvitem{AAAI 2021}{A. Sarkar, \textbf{A. Sarkar}, VN Balasubramanian. ``\textbf{Enhanced Regularizers for Attributional Robustness}''. Association for the Advancement of Artificial Intelligence. \textit{[22 citations]}}

\vspace{0.1cm}

\cvitem{ICML 2019}{A. Chattopadhyay, P. Manupriya, \textbf{A. Sarkar}, VN Balasubramanian. ``\textbf{Neural Network Attributions: A Causal Perspective}''. \textit{Long Oral Presentation} $\cdot$ \textit{[220 citations]}}

\vspace{0.1cm}

\cvitem{WACV 2018}{\textbf{A. Sarkar}, A. Chattopadhyay, P. Howlader, VN Balasubramanian. ``\textbf{Grad-CAM++: Generalized Gradient-based Visual Explanations for Deep Convolutional Networks}''. \\ \textit{\textbf{[4992 citations]} - Highly cited paper}}

\vspace{0.1cm}

\cvitem{Preprint}{\textbf{A. Sarkar}, M. Groth, I. Mason, T. Sasaki, X. Boix. ``\textbf{Deephys: Deep Electrophysiology - Debugging Neural Networks under Distribution Shifts}''.}
% \href{https://arxiv.org/pdf/2303.11912.pdf}{\faExternalLink} $\cdot$ \href{https://deephys.org/}{\textcolor{blue}{Project Website}}}

\vspace{0.2cm}

\subsection{Workshop Papers}
\vspace{0.2cm}

\cvitem{ICML 2026}{\textbf{A. Sarkar}, A. Duran, P. Koo. ``GPA: Generative Population Annealing for Test-Time Sequence Design''. Workshop on Generative \& Agentic AI for Biology}

\cvitem{ICLR 2025}{\textbf{A. Sarkar}, Y. Kang, N. Somia, P. Koo. ``Understanding DNA Discrete Diffusion for Engineering Regulatory DNA Sequences''. Workshop on AI for Nucleic Acids}

\cvitem{NeurIPS 2024}{\textbf{A. Sarkar}, Z. Tang, C. Zhao, P. Koo. ``Designing DNA With Tunable Regulatory Activity Using Discrete Diffusion''. Workshop on AI for New Drug Modalities}

\cvitem{ICCV 2021}{\textbf{A. Sarkar}, A. Sarkar, VN Balasubramanian. ``Leveraging Test-Time Consensus Prediction for Robustness against Unseen Noise''. Workshop on Adversarial Robustness In The Real World}

\vspace{0.25cm}

% ==================== FEATURED RESEARCH PROJECTS ====================
\section{Featured Research Projects}

\cventry{2025--Present}{\textbf{Generative Population Annealing for Test-Time Sequence Design}}{}{}{}{
Inference-time control of generative models to design sequences beyond the training activity range\\
$\bullet$ Model-agnostic method that steers frozen generators without retraining or model-specific hooks\\
$\bullet$ Population annealing-based test-time sampling that pushes generation past activities seen during training\\
$\bullet$ Accepted at GenBio 2026 (ICML Workshop on Generative \& Agentic AI for Biology)
}

\vspace{0.2cm}

\cventry{2023--Present}{\textbf{DNA Sequence Design with Discrete Diffusion Models}}{}{}{}{
Developing generative models for designing regulatory DNA sequences with tunable activity\\
$\bullet$ Novel application of discrete diffusion models to biological sequence design\\
$\bullet$ Enables controllable generation of DNA sequences with specific regulatory properties\\
$\bullet$ Published at MLCB 2024 (Long Oral) and featured in ICLR 2025 and NeurIPS 2024 workshops
}

\vspace{0.2cm}

\cventry{2022--2023}{\textbf{Deephys: Deep Electrophysiology}}{}{}{}{
Analyzing artificial neuron behaviors under distribution shifts inspired by neuroscience\\
$\bullet$ Interactive visualization tool for debugging deep models\\
% : \href{https://deephys.org/}{\textcolor{blue}{deephys.org}}\\
$\bullet$ Novel framework connecting neuroscience principles to deep learning interpretability\\
$\bullet$ Enables understanding of model failures through neuron-level analysis
}

\vspace{0.2cm}

\cventry{2021--2022}{\textbf{Ante-hoc Explainable Models via Concepts}}{IBM Research Collaboration}{}{}{
Learning to predict and explain jointly through concepts during training\\
$\bullet$ Supports multiple supervision levels: unsupervised, weakly-supervised, and fully-supervised\\
$\bullet$ Semantic segmentation of latent concepts for human-understandable explanations\\
$\bullet$ Achieves significant performance improvements while maintaining interpretability
}

\vspace{0.2cm}

\cventry{2020--2021}{\textbf{Teacher-Guided Adversarial Robustness}}{}{}{}{
Non-iterative robust training method without generating adversarial examples\\
$\bullet$ 10x faster training compared to traditional adversarial training methods\\
$\bullet$ Maintains high natural accuracy (minimal trade-off) while improving robustness\\
$\bullet$ Enforces attribution alignment to actual objects, reducing attack surface
}

\vspace{0.2cm}

\cventry{2019--2021}{\textbf{Enhanced Attributional Robustness}}{}{}{}{
Novel regularizers for training models with robust attribution maps\\
$\bullet$ Attribution-based contrastive regularizer enforcing realistic pixel importance distributions\\
$\bullet$ Attribution change-based regularizer for stability under imperceptible perturbations\\
$\bullet$ Achieved state-of-the-art results across multiple robustness benchmarks
}

\vspace{0.2cm}

\cventry{2018--2019}{\textbf{Causal Attribution Methods}}{}{}{}{
Interpreting neural networks through Structural Causal Models (SCMs)\\
$\bullet$ Efficient calculation of interventional expectations and causal attributions\\
$\bullet$ Provides both global understanding and local justifications of model decisions\\
$\bullet$ ICML 2019 Long Oral presentation with 200+ citations
}

\vspace{0.2cm}

\cventry{2017--2018}{\textbf{Grad-CAM++: Visual Explanations for Deep Networks}}{}{}{}{
Generalized gradient-based visual explanation method for CNN architectures\\
$\bullet$ Improved localization and faithfulness compared to existing CAM-based methods\\
$\bullet$ Handles multiple objects and better coverage of object regions\\
$\bullet$ Most cited work with $\sim$5000 citations, widely adopted in computer vision community
}

\vspace{0.25cm}

% ==================== VISITING POSITIONS ====================
\section{Visiting Positions}

\cventry{Jun--Jul 2017}{Visiting Research Student}{University of Tokyo}{Japan}{}{
\textbf{Machine Intelligence Lab} $\cdot$ Advisor: Prof. Tatsuya Harada\\
$\bullet$ Sakura Science Plan Award recipient\\
$\bullet$ Explored causal inference applications in machine learning for model explainability\\
$\bullet$ Developed foundational understanding of causality in deep learning interpretability
}

\vspace{0.2cm}

% ==================== INDUSTRY EXPERIENCE ====================
\section{Industry Experience}

\cventry{May--Jul 2019}{Research Intern}{IBM India Research Lab}{Bangalore, India}{}{
\textit{Project:} Self-explaining neural networks with meaningful concepts\\
$\bullet$ Developed interpretable AI systems with human-understandable building blocks\\
$\bullet$ Led to CVPR 2022 publication on ante-hoc explainable models
}

\cventry{Jun 2010--Dec 2012}{System Engineer}{IBM India Pvt. Ltd.}{Kolkata, India}{}{
Technical consultant for Oracle Applications 11i and Oracle Financials\\
$\bullet$ Oracle e-Business Suite Certification from Oracle University
}

\vspace{0.2cm}

% ==================== HONORS & AWARDS ====================
\section{Honors \& Awards}

\cvitem{2022}{\textbf{\href{https://ikdd.acm.org/awards-2022.php}{\textcolor{blue}{IKDD Best Doctoral Dissertation Award in Data Science}}} $\cdot$ \textit{Winner}}
\cvitem{2022}{\textbf{NASSCOM AI Gamechanger Award} $\cdot$ AI Research (DL Algorithms) $\cdot$ \textit{Runner-up}}
\cvitem{2022}{CVPR travel Grant}
\cvitem{2021}{Selected for WACV Doctoral Consortium}
\cvitem{2019}{ICML Travel Grant}
\cvitem{2019}{IBM Research Internship on AI Explainability and Active Learning}
\cvitem{2018}{Machine Learning Summer School $\cdot$ Universidad Autónoma de Madrid, Spain}
\cvitem{2017}{Sakura Science Plan Award $\cdot$ University of Tokyo Research Internship}

\vspace{0.2cm}

% ==================== TEACHING EXPERIENCE ====================
\section{Teaching Experience}

\cventry{2016--2022}{Teaching \& Research Assistant}{IIT Hyderabad}{}{}{
$\bullet$ \textbf{Deep Learning for Computer Vision} (NPTEL National Platform) - \href{https://onlinecourses.nptel.ac.in/noc22_cs76/preview}{\textcolor{blue}{Course Link}}\\
$\bullet$ Deep Learning for Vision (CS5370) $\cdot$ 100+ graduate and undergraduate students\\
$\bullet$ Optimization Methods in Machine Learning (CS6230) $\cdot$ 25+ students\\
$\bullet$ Deep Learning (CS5480) $\cdot$ 90+ students\\
$\bullet$ Applied Machine Learning (CS6510) $\cdot$ 120+ students
}

\vspace{0.2cm}

% ==================== SERVICE ====================
\section{Professional Service}

\subsection{Conference Reviewing}
\cvitem{}{ICML (2025), NeurIPS (2022-2026), CVPR (2021-2023), ICLR (2022), ICCV (2021), WACV (2022)}

\subsection{Community Service}
\cvitem{}{ICML 2020 Student Volunteer}
\cvitem{}{ACM India Student Chapter (Former Member)}

\vspace{0.2cm}

% ==================== TECHNICAL SKILLS ====================
\section{Technical Skills}
% \cvdoubleitem{\textbf{Specialties}}{Explainable AI, Generative Models}{\textbf{Domains}}{Computer Vision, Computational Biology}

\cvdoubleitem{\textbf{Languages}}{Python, C/C++, R, MATLAB, SQL/PL-SQL}{\textbf{ML/DL}}{PyTorch, TensorFlow, Keras}
\cvdoubleitem{\textbf{Tools}}{Git, LaTeX, Linux, Jupyter}{\textbf{Libraries}}{NumPy, Pandas, Scikit-learn, Matplotlib}

\vspace{0.2cm}

% ==================== SELECTED TALKS & PRESENTATIONS ====================
\section{Selected Talks \& Presentations}

\cvitem{2025}{``Designing DNA With Tunable Regulatory Activity'' $\cdot$ QB/AI Seminar $\cdot$ Cold Spring Harbor Laboratory}
\cvitem{2024}{``Designing DNA With Tunable Regulatory Activity'' $\cdot$ MLCB 2024 $\cdot$ Long Oral}
\cvitem{2022}{``Deephys: Deep Electrophysiology'' $\cdot$ Internal Workshop $\cdot$ Fujitsu Research, Kawasaki, Japan}
\cvitem{2022}{``IKDD Best Doctoral Dissertation Award Talk'' $\cdot$ Data Science Symposium}
\cvitem{2022}{``Learning Ante-hoc Explainable Models'' $\cdot$ CVPR 2022 $\cdot$ New Orleans, USA}
\cvitem{2021}{``Adversarial Robustness without Adversarial Training'' $\cdot$ NeurIPS 2021 $\cdot$ Virtual}
\cvitem{2021}{``Doctoral Consortium Presentation'' $\cdot$ WACV 2021}
\cvitem{2019}{``Neural Network Attributions: A Causal Perspective'' $\cdot$ ICML 2019 $\cdot$ Long Beach, CA}
\cvitem{2018}{``Grad-CAM++: Visual Explanations for Deep Networks'' $\cdot$ WACV 2018 $\cdot$ Lake Tahoe, NV}

% ==================== REFERENCES ====================
\section{References}

\cvitem{}{Available upon request}

% \vfill
% \begin{center}
% {\footnotesize \textit{Last updated: \today}}
% \end{center}

\end{document}